\scanpage{21}

s'assurer de la connexité dans le cas où $\Lambda=\mathbb{Z}$.

Un complément dans le cas $1/2$ local : \deleted{\unsure{Soit|On}} On sait que \deleted{$x,y$} $\forall x,y\in M^{\unsure{\alpha|*}}$, $\exists z\in M^{\alpha}$, tels que

\[
(x,z)\in\operatorname{Rep}(M),\quad(z,y)\in\operatorname{Rep}(M)
\]

\deleted{\unsure{soit|On}} on aura donc

\[
x\wedge z=\lambda e_M,\quad z\wedge y=\mu e_M\qquad\lambda,\mu\in\Lambda^*
\]

et quitte à remplacer $z$ par $1/\lambda\,z$, \unsure{OPS|OK|on a|\deleted{OPS}}

\[
x\wedge z=e_M,\quad z\wedge y=\mu e_M
\]

et par l'argument précédent on trouvera $z'\in M$ tel qu'on ait aussi $z\wedge z'=z'\wedge y=e_M$\unsure{)}, donc

\[
x\wedge z=z\wedge z'=z'\wedge y=e_M\tag*{(47)}
\]

ce qui montre que deux sommets de $SF_M$ (ou de $F_M$) sont à une distance combinatoire $\leq3$. (NB. La borne est atteinte pour l'icosaèdre, correspondant au cas de \deleted{$\unsure{\simeq|\mathbb{Z}}$} $\Lambda=\mathbb{Z}/5\mathbb{Z}$.)\unsure{.}

Pour le cas de $\Lambda=\mathbb{Z}$, particulièrement intéressant de par son rôle ``universel'', il est

% archive sheet number: 20
